\chapter{Elon Lindenstrauss --- Fields Medal 2010}\label{ch2010elonlindenstrauss:elon-lindenstrauss-fields-medal-2010-a-comprehensive-report-on-his-work-and-impact}

\begin{center}\rule{0.5\linewidth}{0.5pt}\end{center}

\textbf{Prepared:} July 2026 \textbf{Subject:} Elon Lindenstrauss (b.
August 1, 1970, Jerusalem, Israel) \textbf{Position:} Professor,
Einstein Institute of Mathematics, Hebrew University of Jerusalem;
Professor, School of Mathematics, Institute for Advanced Study
\textbf{Doctoral Advisor:} Benjamin Weiss \textbf{Fields Medal
Citation:} \emph{``For his results on measure rigidity in ergodic
theory, and their applications to number theory.''}

\begin{center}\rule{0.5\linewidth}{0.5pt}\end{center}

\section{Table of Contents}\label{ch2010elonlindenstrauss:table-of-contents}

\begin{enumerate}
\def\labelenumi{\arabic{enumi}.}
\tightlist
\item
 Introduction
\item
 Biographical Background and Education
\item
 Ergodic Theory: Foundations and Context
\item
 Measure Rigidity: The Furstenberg--Margulis Conjecture
\item
 The Einsiedler--Katok--Lindenstrauss Theorem
\item
 The Littlewood Conjecture
\item
 Quantum Unique Ergodicity: The Rudnick--Sarnak Conjecture
\item
 Arithmetic Quantum Chaos
\item
 Lindenstrauss's Proof of Arithmetic Quantum Unique Ergodicity
\item
 Entropy Arguments in Ergodic Theory
\item
 The Bourgain--Lindenstrauss Collaboration
\item
 Higher Rank Diagonal Actions
\item
 SL(2,R) Actions on Moduli Space
\item
 Applications to Diophantine Approximation
\item
 Further Number-Theoretic Applications
\item
 Broader Impact and Influence
\item
 Awards and Recognition
\item
 Selected Major Publications
\item
 Ongoing Work and Future Directions
\item
 Conclusion
\item
 References and Further Reading
\end{enumerate}

\begin{prereqs} Ergodic theory (measure-preserving transformations, ergodic theorems), basic number theory (Diophantine approximation). \end{prereqs}

\begin{objectives} Explain the ergodic proof of the arithmetic Bogomolov conjecture; describe the Littlewood conjecture and its resolution; understand the connection between dynamics and number theory. \end{objectives}


\begin{center}\rule{0.5\linewidth}{0.5pt}\end{center}

\section{1. Introduction}\label{ch2010elonlindenstrauss:introduction}

Elon Lindenstrauss was awarded the Fields Medal in 2010 for his
fundamental contributions to ergodic theory and its applications to
number theory. His work is characterized by a rare combination of deep
abstract theory and striking concrete applications: he developed
new tools in ergodic theory, particularly the theory of measure
rigidity for higher-rank abelian group actions, and then applied these
tools to solve long-standing problems in number theory, including the
Littlewood conjecture (with Einsiedler and Katok) and the arithmetic
quantum unique ergodicity conjecture (with Bourgain and later
Soundararajan).

Lindenstrauss is the first Israeli mathematician to win the Fields
Medal. He comes from a notable mathematical family: his father, Joram
Lindenstrauss, was a renowned mathematician whose name appears in the
Johnson--Lindenstrauss lemma, and his sister Ayelet is also a
mathematician.

His work bridges two fields that had previously been considered
distinct. Ergodic theory --- the study of the statistical behavior of
dynamical systems --- had long been recognized as having potential
applications to number theory, but realizing this potential required the
development of new and powerful rigidity theorems. Lindenstrauss
provided these theorems, and in doing so transformed both ergodic theory
and its relationship to number theory.

\begin{center}\rule{0.5\linewidth}{0.5pt}\end{center}

\section{2. Biographical Background and
Education}\label{ch2010elonlindenstrauss:biographical-background-and-education}

\subsection{2.1 Early Life}\label{ch2010elonlindenstrauss:early-life}

Elon Lindenstrauss was born on August 1, 1970, in Jerusalem, Israel,
into a family of mathematicians. His father, Joram Lindenstrauss, was a
professor of mathematics at the Hebrew University of Jerusalem, known
for his work in functional analysis (the Johnson--Lindenstrauss lemma).
His mother, Naomi Lindenstrauss, is a computer scientist. His sister
Ayelet Lindenstrauss is also a mathematician.

Growing up in an intellectually rich environment, Lindenstrauss was
exposed to advanced mathematics from an early age. He attended the
Hebrew University Secondary School and represented Israel at the 1988
International Mathematical Olympiad in Canberra, Australia, where he
earned a bronze medal.

\subsection{2.2 Military Service and University
Education}\label{ch2010elonlindenstrauss:military-service-and-university-education}

Following high school, Lindenstrauss was accepted into the Talpiot
program of the Israel Defense Forces, an elite program for gifted
soldiers that combines military service with advanced academic studies.
He studied at the Hebrew University of Jerusalem, earning a BSc in
Mathematics and Physics in 1991 and a master's degree in mathematics in
1995.

\subsection{2.3 Doctoral Studies}\label{ch2010elonlindenstrauss:doctoral-studies}

Lindenstrauss completed his PhD at the Hebrew University in 1999 under
the supervision of Benjamin Weiss, a leading figure in ergodic theory.
His thesis, ``Entropy Properties of Dynamical Systems,'' already showed
his deep interest in the fundamental structures of ergodic theory.

\subsection{2.4 Academic Career}\label{ch2010elonlindenstrauss:academic-career}

After his PhD, Lindenstrauss held positions at the Institute for
Advanced Study in Princeton (1999--2001), Stanford University (Szego
Assistant Professor, 2001--2003), and the Courant Institute at NYU
(visiting member, 2003--2005). He joined the Princeton University
faculty in 2004 and was simultaneously a Clay Mathematics Institute Long
Term Prize Fellow. In 2008, he also accepted a professorship at the
Einstein Institute of Mathematics at the Hebrew University. In 2024, he
was appointed a permanent faculty member in the School of Mathematics of
the Institute for Advanced Study.

\begin{center}\rule{0.5\linewidth}{0.5pt}\end{center}

\section{3. Ergodic Theory: Foundations and
Context}\label{ch2010elonlindenstrauss:ergodic-theory-foundations-and-context}

\subsection{3.1 What Is Ergodic Theory?}\label{ch2010elonlindenstrauss:what-is-ergodic-theory}

Ergodic theory studies the statistical behavior of measure-preserving
transformations. Given a measure space \((X, \mu)\) and a transformation
\(T: X \to X\) that preserves \(\mu\) (i.e., \(\mu(T^{-1}(A)) = \mu(A)\)
for all measurable sets \(A\)), ergodic theory asks about the behavior
of orbits \(\{T^n(x): n \in \mathbb{N}\}\).

The fundamental theorem of ergodic theory, the Birkhoff ergodic theorem
(1931), states that for an ergodic transformation, time averages
converge to space averages:

\[
\lim_{N \to \infty} \frac{1}{N} \sum_{n=0}^{N-1} f(T^n(x)) = \int_X f \, d\mu
\]

for almost every \(x\) and any integrable function \(f\).

\subsection{3.2 Higher Rank Actions}\label{ch2010elonlindenstrauss:higher-rank-actions}

Most of the early developments in ergodic theory focused on actions of
\(\mathbb{Z}\) (a single transformation) or \(\mathbb{R}\) (a flow). But
many natural systems have symmetries given by higher-rank groups like
\(\mathbb{Z}^k\) or \(\mathbb{R}^k\) with \(k \ge 2\). The study of
higher-rank actions turned out to be qualitatively different from the
rank-one case, exhibiting much stronger rigidity properties.

\subsection{3.3 The Furstenberg
Conjecture}\label{ch2010elonlindenstrauss:the-furstenberg-conjecture}

In the 1960s, Hillel Furstenberg initiated the study of higher-rank
diagonal actions on tori and homogeneous spaces. He conjectured that
certain natural higher-rank actions are ``rigid'' in the sense that any
invariant measure must be of a very special form --- specifically, Haar
measure on a closed subgroup. This idea, later developed by Margulis and
others, became the foundation of the theory of measure rigidity.

\begin{center}\rule{0.5\linewidth}{0.5pt}\end{center}

\section{4. Measure Rigidity: The Furstenberg--Margulis
Conjecture}\label{ch2010elonlindenstrauss:measure-rigidity-the-furstenbergmargulis-conjecture}

\subsection{4.1 The Setting}\label{ch2010elonlindenstrauss:the-setting}

Consider the action of the diagonal subgroup
\(A \subset \mathrm{SL}(n, \mathbb{R})\) on the homogeneous space
\(\mathrm{SL}(n, \mathbb{Z}) \backslash \mathrm{SL}(n, \mathbb{R})\).
For \(n \ge 3\) (so the action has rank at least 2), the diagonal action
is conjectured to be ``rigid'': any \(A\)-invariant and ergodic
probability measure should be algebraic, meaning it is the Haar measure
on a closed subgroup.

This is the Furstenberg--Margulis conjecture on measure rigidity for
higher-rank diagonal actions.

\subsection{4.2 Partial Results}\label{ch2010elonlindenstrauss:partial-results}

Before Lindenstrauss's work, partial results were known: - Furstenberg
(1967) proved rigidity for the \(\times 2, \times 3\) action on the
circle. - Margulis proved that any \(A\)-invariant measure that is also
invariant under a large subgroup must be algebraic. - Ratner (1991)
proved a complete classification for actions of unipotent groups, but
the diagonal (Cartan) case remained open.

\subsection{4.3 Lindenstrauss's
Contribution}\label{ch2010elonlindenstrauss:lindenstrausss-contribution}

Lindenstrauss made fundamental contributions to the measure rigidity
problem. In his 2001 paper on quantum unique ergodicity and in
subsequent work, he developed a powerful entropy-based approach to
measure rigidity. The key idea is to use \textbf{Kolmogorov--Sinai
entropy} to distinguish invariant measures. By showing that any
\(A\)-invariant measure with positive entropy must be algebraic, he
established a partial but extremely useful version of the rigidity
conjecture.

\begin{center}\rule{0.5\linewidth}{0.5pt}\end{center}

\section{5. The Einsiedler--Katok--Lindenstrauss
Theorem}\label{ch2010elonlindenstrauss:the-einsiedlerkatoklindenstrauss-theorem}

\subsection{5.1 The Statement}\label{ch2010elonlindenstrauss:the-statement}

In a landmark 2006 paper, Manfred Einsiedler, Anatole Katok, and Elon
Lindenstrauss proved:

\textbf{Theorem:} Let \(\mu\) be an \(A\)-invariant and ergodic
probability measure on
\(\mathrm{SL}(n, \mathbb{Z}) \backslash \mathrm{SL}(n, \mathbb{R})\) for
\(n \ge 3\). If \(\mu\) has positive entropy with respect to some
element of \(A\), then \(\mu\) is algebraic (i.e., it is the Haar
measure on a closed subgroup of the form
\(L \backslash \mathrm{SL}(n, \mathbb{R})\)).

\subsection{5.2 The Method}\label{ch2010elonlindenstrauss:the-method}

The proof combines: - \textbf{Entropy theory:} Using the relationship
between entropy, conditional measures, and the geometry of the
homogeneous space. - \textbf{Product structure:} Exploiting the fact
that the action of \(A\) decomposes into one-dimensional ``root''
actions. - \textbf{Nonlinear dynamics:} Using the fact that the
horospherical subgroups generate the full group.

\subsection{5.3 Significance}\label{ch2010elonlindenstrauss:significance}

This theorem was the first major result on measure rigidity for diagonal
actions without any additional invariance assumptions (beyond positive
entropy). It showed that the Furstenberg--Margulis conjecture is true
under the mild extra condition of positive entropy --- a condition that
is almost always satisfied in applications.

\begin{center}\rule{0.5\linewidth}{0.5pt}\end{center}

\section{6. The Littlewood Conjecture}\label{ch2010elonlindenstrauss:the-littlewood-conjecture}

\subsection{6.1 The Conjecture}\label{ch2010elonlindenstrauss:the-conjecture}

The Littlewood conjecture in Diophantine approximation states that for
any two real numbers \(\alpha, \beta\):

\[
\liminf_{q \to \infty} q \cdot \|q\alpha\| \cdot \|q\beta\| = 0
\]

where \(\|x\|\) denotes the distance from \(x\) to the nearest integer.
In words: one can simultaneously approximate \(\alpha\) and \(\beta\) by
rationals with the same denominator, with an error product that decays
faster than \(1/q\).

Despite intensive study, the conjecture remains open in full generality.

\subsection{6.2 The Einsiedler--Katok--Lindenstrauss
Result}\label{ch2010elonlindenstrauss:the-einsiedlerkatoklindenstrauss-result}

In the same 2006 paper, Einsiedler, Katok, and Lindenstrauss used their
measure rigidity theorem to prove:

\textbf{Theorem:} The set of pairs \((\alpha, \beta)\) for which the
Littlewood conjecture fails has Hausdorff dimension zero.

In particular, the conjecture holds for almost every pair of real
numbers --- and the exceptional set is extremely small (dimension zero).

\subsection{6.3 Why Measure Rigidity
Applies}\label{ch2010elonlindenstrauss:why-measure-rigidity-applies}

The connection between the Littlewood conjecture and measure rigidity
comes through the space of lattices. The pair \((\alpha, \beta)\)
corresponds to a point in
\(\mathrm{SL}(3, \mathbb{Z}) \backslash \mathrm{SL}(3, \mathbb{R})\)
via:

\[
\begin{pmatrix} 1 & \alpha & \beta \\ 0 & 1 & 0 \\ 0 & 0 & 1 \end{pmatrix}
\]

The failure of the Littlewood conjecture for \((\alpha, \beta)\) implies
that the trajectory of this point under the diagonal action stays in a
compact region, which in turn implies the existence of a nontrivial
invariant measure. The rigidity theorem then forces this measure to be
algebraic, which gives control over the exceptional set.

\subsection{6.4 Impact}\label{ch2010elonlindenstrauss:impact}

This result was immediately recognized as a major breakthrough. It
showed that the Littlewood conjecture --- one of the central open
problems in Diophantine approximation --- is ``almost true'' in a
precise sense, and it demonstrated the power of ergodic theory to
address classical problems in number theory.

\begin{center}\rule{0.5\linewidth}{0.5pt}\end{center}

\section{7. Quantum Unique Ergodicity: The Rudnick--Sarnak
Conjecture}\label{ch2010elonlindenstrauss:quantum-unique-ergodicity-the-rudnicksarnak-conjecture}

\subsection{7.1 Quantum Chaos}\label{ch2010elonlindenstrauss:quantum-chaos}

Quantum chaos studies the behavior of quantum systems whose classical
counterparts are chaotic. A central question is: how do eigenfunctions
of the Laplace operator (the quantum energy states) behave in the
semiclassical limit (high energy)?

\subsection{7.2 The Manifold Setting}\label{ch2010elonlindenstrauss:the-manifold-setting}

For a compact Riemannian manifold \(M\) with ergodic geodesic flow, the
quantum ergodicity theorem (Shnirelman, Zelditch, Colin de Verdière,
1970s--1980s) states that almost all eigenfunctions of the Laplacian
equidistribute in the high-energy limit. However, this theorem allows a
sparse subsequence of eigenfunctions to behave differently
(``scarring'').

The question of whether there are any exceptional subsequences --- the
\textbf{quantum unique ergodicity} (QUE) conjecture --- was posed by
Rudnick and Sarnak in 1994. They conjectured that for manifolds of
negative curvature, \textbf{all} eigenfunctions equidistribute: there
are no scars.

\subsection{7.3 Arithmetic Quantum
Chaos}\label{ch2010elonlindenstrauss:arithmetic-quantum-chaos}

Rudnick and Sarnak also proposed an arithmetic version of the
conjecture. On arithmetic hyperbolic surfaces (quotients of the
hyperbolic plane by congruence subgroups of
\(\mathrm{SL}(2, \mathbb{Z})\)), there is an additional symmetry: the
Hecke operators. These commuting operators allow one to select a basis
of joint eigenfunctions of the Laplacian and the Hecke operators. The
\textbf{arithmetic quantum unique ergodicity} (AQUE) conjecture states
that these Hecke eigenfunctions equidistribute.

\subsection{7.4 The Obstacle}\label{ch2010elonlindenstrauss:the-obstacle}

The general QUE conjecture for negatively curved manifolds remains open.
The arithmetic case seemed slightly more tractable because the Hecke
operators provide additional structure, but the problem was still
extremely difficult.

\begin{center}\rule{0.5\linewidth}{0.5pt}\end{center}

\section{8. Arithmetic Quantum Chaos}\label{ch2010elonlindenstrauss:arithmetic-quantum-chaos-1}

\subsection{8.1 The Modular Surface}\label{ch2010elonlindenstrauss:the-modular-surface}

The primary setting for arithmetic quantum chaos is the modular surface
\(\mathrm{SL}(2, \mathbb{Z}) \backslash \mathbb{H}\), where
\(\mathbb{H}\) is the hyperbolic plane. This non-compact surface has a
continuous spectrum, but one can study either compact surfaces (by
considering congruence subgroups) or the continuous spectrum.

\subsection{8.2 Hecke Eigenfunctions}\label{ch2010elonlindenstrauss:hecke-eigenfunctions}

On arithmetic surfaces, the Hecke operators \(T_n\) form a commutative
algebra of self-adjoint operators that commute with the Laplacian. A
Hecke-Maaß cusp form is a joint eigenfunction of the Laplacian and all
Hecke operators. These are the fundamental objects of study in
arithmetic quantum chaos.

\subsection{8.3 Prior Results}\label{ch2010elonlindenstrauss:prior-results}

Before Lindenstrauss's work, the following was known: - Quantum
ergodicity (QE) holds for the modular surface (Zelditch, Luo--Sarnak). -
The QUE conjecture for Eisenstein series was proved by Luo--Sarnak. -
Nalini Anantharaman had proved that any quantum limit must have positive
entropy.

\begin{center}\rule{0.5\linewidth}{0.5pt}\end{center}

\section{9. Lindenstrauss's Proof of Arithmetic Quantum Unique
Ergodicity}\label{ch2010elonlindenstrauss:lindenstrausss-proof-of-arithmetic-quantum-unique-ergodicity}

\subsection{9.1 The Main Theorem}\label{ch2010elonlindenstrauss:the-main-theorem}

In a 2006 Annals of Mathematics paper, Lindenstrauss proved:

\textbf{Theorem:} Any quantum limit of Hecke-Maaß cusp forms on the
modular surface is a constant multiple of the hyperbolic volume form.

In particular, he proved the arithmetic QUE conjecture for the full
modular surface \(\mathrm{SL}(2, \mathbb{Z}) \backslash \mathbb{H}\).

\subsection{9.2 The Strategy}\label{ch2010elonlindenstrauss:the-strategy}

The proof has three main components:

\begin{enumerate}
\def\labelenumi{\arabic{enumi}.}
\item
 \textbf{Entropy bounds:} Using the Hecke operators, Lindenstrauss
 proved that any quantum limit measure must have positive entropy. This
 is the most innovative part of the proof, using the commutativity of
 Hecke operators to force the eigenfunctions to spread out.
\item
 \textbf{Measure rigidity:} The positive entropy quantum limit measure
 is invariant under the diagonal action (the geodesic flow). By
 applying the Einsiedler--Katok rigidity theorem, Lindenstrauss deduced
 that the measure must be algebraic.
\item
 \textbf{Classification:} The only algebraic invariant measures on the
 modular surface that could arise as quantum limits are the hyperbolic
 volume form.
\end{enumerate}

\subsection{9.3 The Entropy Argument}\label{ch2010elonlindenstrauss:the-entropy-argument}

The key insight is that the Hecke operators act as ``correspondences''
on the surface, and that the commutativity of the Laplacian and the
Hecke operators forces any quantum limit to have certain invariance
properties. Lindenstrauss used this to show that any quantum limit has
positive entropy, which then triggers the rigidity theorem.

\subsection{9.4 Completion by
Soundararajan}\label{ch2010elonlindenstrauss:completion-by-soundararajan}

For compact arithmetic surfaces (congruence surfaces of higher level),
the QUE conjecture was completed by Kannan Soundararajan in 2010, using
a different method (subconvexity bounds for L-functions).
Lindenstrauss's approach already handled the ``non-degenerate'' case;
Soundararajan filled in the remaining gaps.

\begin{center}\rule{0.5\linewidth}{0.5pt}\end{center}

\section{10. Entropy Arguments in Ergodic
Theory}\label{ch2010elonlindenstrauss:entropy-arguments-in-ergodic-theory}

\subsection{10.1 Kolmogorov--Sinai
Entropy}\label{ch2010elonlindenstrauss:kolmogorovsinai-entropy}

Entropy is a measure of the ``randomness'' or ``complexity'' of a
dynamical system. For a measure-preserving transformation \(T\) on a
probability space \((X, \mu)\), the Kolmogorov--Sinai entropy
\(h_\mu(T)\) measures the average information gained per iteration of
the transformation.

In a fundamental sense, entropy distinguishes ``chaotic'' systems
(positive entropy) from ``predictable'' ones (zero entropy). The
geodesic flow on a compact hyperbolic surface has positive entropy with
respect to the Liouville measure.

\subsection{10.2 Lindenstrauss's Entropy
Techniques}\label{ch2010elonlindenstrauss:lindenstrausss-entropy-techniques}

Lindenstrauss developed novel entropy-based techniques for studying
invariant measures:

\begin{itemize}
\tightlist
\item
 \textbf{Entropy of quantum limits:} Showing that any limit of
 eigenfunctions must have positive entropy when the eigenfunctions are
 Hecke eigenfunctions.
\item
 \textbf{Entropy and rigidity:} Using entropy to trigger measure
 rigidity theorems.
\item
 \textbf{Conditional entropy:} Using the relationship between entropy
 and conditional measures on invariant foliations.
\end{itemize}

These techniques have become standard tools in ergodic theory and have
been applied to a wide range of problems beyond their original context.

\begin{center}\rule{0.5\linewidth}{0.5pt}\end{center}

\section{11. The Bourgain--Lindenstrauss
Collaboration}\label{ch2010elonlindenstrauss:the-bourgainlindenstrauss-collaboration}

\subsection{11.1 Joint Work on QUE}\label{ch2010elonlindenstrauss:joint-work-on-que}

In a series of papers with Jean Bourgain, Lindenstrauss extended his QUE
results to more general settings: - \textbf{Compact congruence
surfaces:} Handling the case of compact arithmetic surfaces. -
\textbf{Higher-dimensional manifolds:} Generalizations to
\(\mathbb{H}^d\) and other symmetric spaces. - \textbf{Division
algebras:} QUE for quotients of division algebras.

The Bourgain--Lindenstrauss papers developed new entropy bounds that
applied to these more general settings.

\subsection{11.2 The Entropy Bound}\label{ch2010elonlindenstrauss:the-entropy-bound}

The main technical achievement was an entropy bound for joint
eigenfunctions of the Laplacian and Hecke operators on congruence
surfaces. Bourgain and Lindenstrauss proved:

\[
h(\mu) \ge c > 0
\]

for any quantum limit measure \(\mu\), where \(c\) is an absolute
constant independent of the sequence of eigenfunctions.

\subsection{11.3 Impact}\label{ch2010elonlindenstrauss:impact-1}

These results established arithmetic QUE in the broadest possible
setting and demonstrated the power of combining analytic number theory
(Bourgain's expertise) with ergodic theory (Lindenstrauss's expertise).

\begin{center}\rule{0.5\linewidth}{0.5pt}\end{center}

\section{12. Higher Rank Diagonal
Actions}\label{ch2010elonlindenstrauss:higher-rank-diagonal-actions}

\subsection{12.1 The General Theory}\label{ch2010elonlindenstrauss:the-general-theory}

Beyond the specific applications, Lindenstrauss has developed a general
theory of higher-rank diagonal actions on homogeneous spaces. This
theory includes:

\begin{itemize}
\tightlist
\item
 \textbf{Classification of invariant measures:} Under what conditions
 must an invariant measure be algebraic?
\item
 \textbf{Entropy and rigidity:} The relationship between positive
 entropy and algebraic structure.
\item
 \textbf{Mixings and joinings:} Spectral properties of higher-rank
 actions.
\end{itemize}

\subsection{12.2 Key Results}\label{ch2010elonlindenstrauss:key-results}

Some of the key results include: - If a measure is invariant under a
higher-rank diagonal action and has positive entropy for some element,
it is algebraic (with Einsiedler and Katok). - Any measure invariant
under the full diagonal group is a combination of algebraic measures. -
New proofs and generalizations of Ratner's theorem using entropy
methods.

\subsection{12.3 The Furstenberg--Margulis
Conjecture}\label{ch2010elonlindenstrauss:the-furstenbergmargulis-conjecture}

The full Furstenberg--Margulis conjecture --- that any \(A\)-invariant
measure on
\(\mathrm{SL}(n, \mathbb{Z}) \backslash \mathrm{SL}(n, \mathbb{R})\) is
algebraic, without the positive entropy hypothesis --- remains open.
Lindenstrauss's work reduces the problem to the zero-entropy case, which
is still mysterious.

\begin{center}\rule{0.5\linewidth}{0.5pt}\end{center}

\section{13. SL(2,R) Actions on Moduli
Space}\label{ch2010elonlindenstrauss:sl2r-actions-on-moduli-space}

\subsection{13.1 The Moduli Space of Abelian
Differentials}\label{ch2010elonlindenstrauss:the-moduli-space-of-abelian-differentials}

Lindenstrauss has also made contributions to the study of the SL(2,R)
action on the moduli space of abelian differentials --- the same setting
in which Mirzakhani and Eskin proved the Magic Wand Theorem.

\subsection{13.2 Applications of Measure
Rigidity}\label{ch2010elonlindenstrauss:applications-of-measure-rigidity}

Using his measure rigidity techniques, Lindenstrauss (together with
co-authors) has studied: - Invariant measures for the SL(2,R) action on
moduli space. - Classification of orbit closures in certain cases. -
Connections to billiards and interval exchange transformations.

This work shows the wide applicability of the entropy-rigidity framework
beyond the original setting of homogeneous spaces.

\begin{center}\rule{0.5\linewidth}{0.5pt}\end{center}

\section{14. Applications to Diophantine
Approximation}\label{ch2010elonlindenstrauss:applications-to-diophantine-approximation}

\subsection{14.1 Beyond Littlewood}\label{ch2010elonlindenstrauss:beyond-littlewood}

The Einsiedler--Katok--Lindenstrauss theorem on the Littlewood
conjecture opened up a new direction in Diophantine approximation: using
ergodic theory to study multidimensional approximation problems.

\subsection{14.2 Related Problems}\label{ch2010elonlindenstrauss:related-problems}

Subsequent work has applied similar techniques to: -
\textbf{Simultaneous approximation:} Approximation of vectors by
rationals. - \textbf{Mixed approximation:} Different rates of
approximation for different coordinates. - \textbf{Nonlinear
approximation:} Approximation by algebraic numbers.

Each of these problems can be translated into a problem about invariant
measures on homogeneous spaces, which can then be attacked using
entropy-rigidity methods.

\subsection{14.3 Effective Bounds}\label{ch2010elonlindenstrauss:effective-bounds}

One limitation of the ergodic-theoretic approach is that it is typically
ineffective: it proves existence but does not give explicit bounds.
Making these results effective --- finding explicit constants in the
Littlewood problem, for example --- is an ongoing challenge.

\begin{center}\rule{0.5\linewidth}{0.5pt}\end{center}

\section{15. Further Number-Theoretic
Applications}\label{ch2010elonlindenstrauss:further-number-theoretic-applications}

\subsection{15.1 Uniform Distribution}\label{ch2010elonlindenstrauss:uniform-distribution}

Lindenstrauss's work on measure rigidity has implications for uniform
distribution problems in number theory. For example: - Distribution of
lattice points on expanding horospheres - Equidistribution of Hecke
points - Distribution of periodic orbits of the geodesic flow

\subsection{15.2 The Sarnak Conjecture}\label{ch2010elonlindenstrauss:the-sarnak-conjecture}

The Sarnak conjecture on the Möbius function --- that the Möbius
function does not correlate with any deterministic sequence --- has
connections to ergodic theory. Lindenstrauss's entropy methods have been
applied to study this conjecture in certain settings.

\subsection{15.3 Spectral Theory}\label{ch2010elonlindenstrauss:spectral-theory}

The QUE result has implications for the spectral theory of automorphic
forms: - Distribution of Maaß cusp forms - Limits of Hecke-Maaß forms -
Connections to the Riemann hypothesis through the Selberg trace formula

\begin{center}\rule{0.5\linewidth}{0.5pt}\end{center}

\section{16. Broader Impact and
Influence}\label{ch2010elonlindenstrauss:broader-impact-and-influence}

\subsection{16.1 Creation of a New Field}\label{ch2010elonlindenstrauss:creation-of-a-new-field}

Lindenstrauss's work has created a new field at the intersection of
ergodic theory and number theory. Before his work, applications of
ergodic theory to number theory were limited to specific problems (e.g.,
the Hardy--Littlewood circle method). His demonstration that deep
rigidity theorems in ergodic theory could solve central problems in
number theory transformed the relationship between the two fields.

\subsection{16.2 Influence on Young
Researchers}\label{ch2010elonlindenstrauss:influence-on-young-researchers}

Lindenstrauss has supervised numerous PhD students and postdocs who have
gone on to become leaders in the field. His entropy-rigidity framework
has been taken up by researchers around the world and applied to
problems ranging from Diophantine approximation to mathematical physics.

\subsection{16.3 The Einsiedler--Katok--Lindenstrauss
Collaboration}\label{ch2010elonlindenstrauss:the-einsiedlerkatoklindenstrauss-collaboration}

The collaboration between Einsiedler, Katok, and Lindenstrauss produced
one of the influential papers in 21st-century dynamics. The
combination of Einsiedler's analytic skills, Katok's deep knowledge of
dynamical systems, and Lindenstrauss's entropy techniques proved
remarkably effective.

\subsection{16.4 The Lindenstrauss Mathematical
Dynasty}\label{ch2010elonlindenstrauss:the-lindenstrauss-mathematical-dynasty}

Lindenstrauss comes from a notable mathematical family: - His father,
Joram Lindenstrauss (d.~2012), was a leading functional analyst. - His
sister, Ayelet Lindenstrauss, is a mathematician specializing in
algebraic topology. - His doctoral advisor, Benjamin Weiss, is one of
the leading figures in ergodic theory. His own students include some of
the most promising young mathematicians in the field.

\begin{center}\rule{0.5\linewidth}{0.5pt}\end{center}

\section{17. Awards and Recognition}\label{ch2010elonlindenstrauss:awards-and-recognition}

\textbf{2010:} - Fields Medal (International Mathematical Union)

\textbf{2009:} - Fermat Prize (Institut de Mathématiques de Toulouse) -
Erdos Prize (Israel Mathematical Union)

\textbf{2008:} - Michael Bruno Memorial Award

\textbf{2004:} - European Mathematical Society Prize

\textbf{2003:} - Clay Mathematics Institute Long Term Prize Fellowship -
Salem Prize

\textbf{2001:} - Blumenthal Award for the Advancement of Research in
Pure Mathematics (American Mathematical Society)

\begin{center}\rule{0.5\linewidth}{0.5pt}\end{center}

\section{18. Selected Major
Publications}\label{ch2010elonlindenstrauss:selected-major-publications}

\begin{enumerate}
\def\labelenumi{\arabic{enumi}.}
\item
 \textbf{``Invariant measures and arithmetic quantum unique
 ergodicity.''} \emph{Annals of Mathematics} 163(1), 165--219, 2006.
 --- \emph{Proof of arithmetic QUE for the full modular surface.}
\item
 \textbf{``Invariant measures for higher-rank hyperbolic abelian
 actions''} (with M. Einsiedler and A. Katok). \emph{Annals of
 Mathematics} 164(2), 513--560, 2006. --- \emph{Measure rigidity for
 diagonal actions and the Littlewood conjecture result.}
\item
 \textbf{``On the quantum unique ergodicity conjecture for
 \(\Gamma \backslash \mathbb{H} \times \mathbb{H}\).''}
 \emph{International Mathematics Research Notices} 2001(17), 2001. ---
 \emph{QUE for products of hyperbolic planes.}
\item
 \textbf{``Entropy in quantum ergodicity''} (with J. Bourgain).
 \emph{Communications in Mathematical Physics} 253(2), 453--469, 2005.
 --- \emph{Entropy bounds for quantum limits.}
\item
 \textbf{``Graph eigenfunctions and quantum unique ergodicity''} (with
 S. Brooks). arXiv:1006.3583, 2010. --- \emph{QUE for congruence
 surfaces.}
\item
 \textbf{``Pointwise estimates for relative entropy''} (with M.
 Einsiedler). \emph{Journal of Modern Dynamics} 1(3), 359--385, 2007.
 --- \emph{Technical foundations for entropy-rigidity.}
\item
 \textbf{``Diagonal actions and positive entropy''} (with M.
 Einsiedler). \emph{Journal of Modern Dynamics} 1(4), 2007. ---
 \emph{General theory of higher-rank diagonal actions.}
\item
 \textbf{``On the joint distribution of the fractional parts of
 \(n\alpha, n\beta\).''} \emph{Journal d'Analyse Mathématique} 90,
 135--155, 2003. --- \emph{Application of ergodic theory to Diophantine
 approximation.}
\item
 \textbf{``The distribution of periodic torus orbits on homogeneous
 spaces''} (with M. Einsiedler and U. Shapira). \emph{Duke Mathematical
 Journal} 160(2), 2007. --- \emph{Counting and equidistribution
 results.}
\end{enumerate}

\begin{center}\rule{0.5\linewidth}{0.5pt}\end{center}

\section{19. Ongoing Work and Future
Directions}\label{ch2010elonlindenstrauss:ongoing-work-and-future-directions}

\subsection{19.1 The Full Furstenberg--Margulis
Conjecture}\label{ch2010elonlindenstrauss:the-full-furstenbergmargulis-conjecture}

The conjecture that all invariant measures for higher-rank diagonal
actions are algebraic, without any entropy assumption, remains open.
Lindenstrauss continues to work on this problem, which is one of the
central questions in ergodic theory.

\subsection{19.2 Effective Results in Diophantine
Approximation}\label{ch2010elonlindenstrauss:effective-results-in-diophantine-approximation}

Making the Littlewood conjecture result effective --- finding explicit
rates of approximation --- is an important open problem. Lindenstrauss's
methods are generally ineffective, but there is ongoing work to combine
them with analytic number theory to obtain effective bounds.

\subsection{19.3 General QUE on Negatively Curved
Manifolds}\label{ch2010elonlindenstrauss:general-que-on-negatively-curved-manifolds}

The general QUE conjecture (without the Hecke operators) remains
completely open. Lindenstrauss's entropy methods provide one approach,
but the general case requires new ideas.

\subsection{19.4 Sparse Equidistribution}\label{ch2010elonlindenstrauss:sparse-equidistribution}

Lindenstrauss has recently been working on problems of sparse
equidistribution --- understanding how sequences become equidistributed
when the number of points grows more slowly than the volume.

\begin{center}\rule{0.5\linewidth}{0.5pt}\end{center}

\section{20. Conclusion}\label{ch2010elonlindenstrauss:conclusion}

Elon Lindenstrauss is one of the influential mathematicians working
at the intersection of ergodic theory and number theory. In a career
spanning three decades, he has:

\begin{itemize}
\tightlist
\item
 Developed the theory of measure rigidity for higher-rank diagonal
 actions
\item
 Proved (with Einsiedler and Katok) that the Littlewood conjecture
 holds for almost all pairs of real numbers
\item
 Proved the arithmetic quantum unique ergodicity conjecture for the
 modular surface
\item
 Created a general framework connecting entropy to rigidity in
 dynamical systems
\item
 Applied ergodic theory to a wide range of problems in number theory
\end{itemize}

The Fields Medal citation --- \emph{``for his results on measure
rigidity in ergodic theory, and their applications to number theory''}
--- captures the essence of his contributions. His work has transformed
both ergodic theory (by providing a powerful new set of tools for
studying invariant measures) and number theory (by demonstrating that
these tools can solve classical problems). The field he has created ---
arithmetic dynamics --- continues to flourish, driven by the methods and
insights he pioneered.

\begin{center}\rule{0.5\linewidth}{0.5pt}\end{center}

\begin{center}\rule{0.5\linewidth}{0.5pt}\end{center}

\section*{Exercises}
\addcontentsline{toc}{section}{Exercises}

\begin{exercise} Prove the Birkhoff ergodic theorem: for a measure-preserving transformation $T$ and $f \\in L^1$, $\\frac{1}{n} \\sum_{k=0}^{n-1} f(T^k x) \\to \\int f \\, d\\mu$ a.e. \end{exercise}

\begin{exercise} Show that the geodesic flow on a hyperbolic surface is ergodic. \end{exercise}

\begin{exercise} Verify that the Littlewood conjecture $\\liminf_{n \\to \\infty} n \\|n\\alpha\\| \\|n\\beta\\| = 0$ holds for almost every pair $(\\alpha, \\beta)$. \end{exercise}

\begin{exercise} Compute the entropy of the geodesic flow on a compact surface of constant negative curvature. \end{exercise}

\begin{center}\rule{0.5\linewidth}{0.5pt}\end{center}

\section*{Summary}
\addcontentsline{toc}{section}{Summary}

Lindenstrauss made fundamental contributions to ergodic theory and its applications to number theory. His ergodic proof of the arithmetic Bogomolov conjecture and his work on the Littlewood conjecture revealed deep connections between dynamical systems and Diophantine approximation, and his development of measure rigidity advanced the field of homogeneous dynamics.

\section{21. References and Further
Reading}\label{ch2010elonlindenstrauss:references-and-further-reading}

\subsection{Official Sources}\label{ch2010elonlindenstrauss:official-sources}

\begin{itemize}
\tightlist
\item
 ICM 2010 Fields Medal Announcement:
 https://www.icm2010.in/prize-winners-2010/fields-medal-elon-lindenstrauss
\item
 IMU Fields Medal Citations:
 https://www.mathunion.org/imu-awards/fields-medal/fields-medal-2010
\end{itemize}

\subsection{Biographical and Overview
Articles}\label{ch2010elonlindenstrauss:biographical-and-overview-articles}

\begin{itemize}
\tightlist
\item
 Princeton University News. ``Lindenstrauss Wins Prestigious Fields
 Medal.'' August 20, 2010.
\item
 Sarnak, P. ``The Work of Elon Lindenstrauss.'' ICM 2010 Laudatio.
\item
 Klarreich, E. ``A Bridge Between the Worlds of Dynamics and Number
 Theory.'' \emph{Quanta Magazine}, 2010.
\end{itemize}

\subsection{Key Technical References}\label{ch2010elonlindenstrauss:key-technical-references}

\begin{itemize}
\tightlist
\item
 Lindenstrauss, E. ``Invariant measures and arithmetic quantum unique
 ergodicity.'' \emph{Annals of Mathematics} 163 (2006), 165--219.
\item
 Einsiedler, M., Katok, A., and Lindenstrauss, E. ``Invariant measures
 for higher-rank hyperbolic abelian actions.'' \emph{Annals of
 Mathematics} 164 (2006), 513--560.
\item
 Bourgain, J. and Lindenstrauss, E. ``Entropy in quantum ergodicity.''
 \emph{Communications in Mathematical Physics} 253 (2005), 453--469.
\item
 Lindenstrauss, E. ``On quantum unique ergodicity for
 \(\Gamma \backslash \mathbb{H} \times \mathbb{H}\).''
 \emph{International Mathematics Research Notices} 2001.
\item
 Einsiedler, M. and Lindenstrauss, E. ``Diagonal actions and positive
 entropy.'' \emph{Journal of Modern Dynamics} 1 (2007).
\end{itemize}

\begin{center}\rule{0.5\linewidth}{0.5pt}\end{center}

\section*{See Also}
\addcontentsline{toc}{section}{See Also}

This chapter relates to the following work by other Fields Medalists:
\begin{itemize}
  \item \textbf{\hyperref[ch1990edwardwitten:edward-witten-fields-medal-1990-a-comprehensive-report]{See Chapter on related work}} --- Lindenstrauss's work on ergodic theory relates.
  \item \textbf{\hyperref[ch2006terencetao:terence-tao-fields-medal-2006-a-comprehensive-report-on-his-work-and-impact]{See Chapter on related work}} --- Tao's work on ergodic theory relates.
\end{itemize}

